\section{Method}\label{sec:method}

Figure~\ref{fig:pipeline} gives the end-to-end pipeline: a class-weighted
out-of-bag forest ensemble supplies the model-derived signals; misclassified
instances are triaged into three categories by out-of-bag confidence and
local-representation rules whose interpretation is grounded in, and validated
against, an aleatoric--epistemic uncertainty decomposition
(\S\ref{sec:method:uncert}--\ref{sec:method:triage});
the misclassified \emph{minority} is then mapped to remedies by the remedy
decomposition (\S\ref{sec:method:remedy}); and the resulting diagnosis is
evaluated against a 16-strategy menu (\S\ref{sec:method:menu},
\S\ref{sec:design}). Table~\ref{tab:algorithm} states the procedure as
pseudocode.

\begin{figure}[!tp]
\centering
\begin{tikzpicture}[
  font=\footnotesize,
  stage/.style={draw, rounded corners, align=center, inner sep=3pt,
                minimum height=8mm, fill=black!4},
  sig/.style={stage, fill=blue!6},
  result/.style={stage, fill=green!8},
  >={Stealth[length=2mm]}, node distance=5mm and 6mm
]
\node[stage, text width=24mm] (data) {Training split\\(stratified, $\le$10k)};
\node[stage, right=of data, text width=32mm] (ens)
  {Class-weighted ensemble\\$M{=}5$ RF $\times$ $T{=}100$\\(category signals OOB)};
\node[stage, right=of ens, text width=24mm] (err)
  {OOB error set\\$E=\{\hat{y}^{\mathrm{OOB}}_i \ne y_i\}$};
\node[sig, below=of ens, text width=92mm] (sigs)
  {Per-error OOB signals: true-class probability (TCP);\quad
   aleatoric/epistemic uncertainty (Eq.~\ref{eq:uncertainty});\quad
   local class ratio};
\node[sig, below=of sigs, text width=92mm] (triage)
  {Triage: \ Cat\,1 noise ($\mathrm{TCP}<\tau_{\text{noise}}$) \ $\cdot$ \
   Cat\,2 data-limited (low class ratio) \ $\cdot$ \
   Cat\,3 boundary/overlap};
\node[sig, below=of triage, xshift=-22mm, text width=48mm] (remedy)
  {Remedy decomposition of \emph{minority} errors\\
   (threshold gate at OOF $\tau^\star$, then triage; Figure~\ref{fig:remedy})};
\node[result, right=of remedy, text width=32mm] (outp)
  {Diagnosis\\(remedy shares)};
\node[result, below=of remedy, xshift=22mm, text width=92mm] (eval)
  {Menu evaluation (\S\ref{sec:design}), $5\times5$ stratified CV: strategy +
   threshold fit on \emph{training} folds $\to$ learner fit $\to$ \emph{held-out}
   test predictions $\to$ default/tuned decisions $\to$ metrics};
\draw[->] (data) -- (ens);
\draw[->] (ens) -- (err);
\draw[->] (err.south) -- (err.south |- sigs.north);
\draw[->] (sigs) -- (triage);
\draw[->] (triage.south -| remedy.north) -- (remedy.north);
\draw[->] (remedy) -- (outp);
\draw[->] (outp.south) -- (outp.south |- eval.north);
\end{tikzpicture}
\caption{End-to-end pipeline. Every signal that drives the categorization (error
set, true-class probability, noise consensus) comes from out-of-bag predictions
of a class-weighted random-forest ensemble (the uncertainty values of
Eq.~\eqref{eq:uncertainty} interpret the categories and are recomputed fully
held-out in \appref{app:oof}); errors are triaged into three categories;
the misclassified minority is mapped to remedies via the threshold-recoverability
gate of \S\ref{sec:method:remedy}. The diagnosis is then evaluated against the
16-strategy menu under the protocol of \S\ref{sec:design}: $5\times5$ stratified
cross-validation in which each strategy (resampling, weighting, or
threshold-moving) and any tuned threshold are fit on the training folds only,
and all reported metrics come from the held-out test folds.}
\label{fig:pipeline}
\end{figure}

\subsection{Uncertainty decomposition via random forest ensembles}\label{sec:method:uncert}

We train an ensemble of $M=5$ random forests \cite{breiman2001random}, each with
$T=100$ trees and minimum leaf size $5$, on the full training set. For each
instance $\mathbf{x}_i$, each tree $t$ in forest $m$ produces a class probability
vector $\hat{p}_{m,t}(\mathbf{x}_i)$. Following Shaker and H\"ullermeier
\cite{shaker2020aleatoric,hullermeier2021aleatoric}, we decompose the total
predictive uncertainty:
\begin{equation}
\underbrace{H[\bar{p}(\mathbf{x}_i)]}_{\text{total}}
= \underbrace{\frac{1}{MT}\sum_{m,t} H[\hat{p}_{m,t}(\mathbf{x}_i)]}_{\text{aleatoric}}
+ \underbrace{H[\bar{p}(\mathbf{x}_i)] - \frac{1}{MT}\sum_{m,t} H[\hat{p}_{m,t}(\mathbf{x}_i)]}_{\text{epistemic}}
\label{eq:uncertainty}
\end{equation}
where $\bar{p}(\mathbf{x}_i) = \frac{1}{MT}\sum_{m,t}\hat{p}_{m,t}(\mathbf{x}_i)$
is the ensemble-averaged prediction and $H[\cdot]$ denotes Shannon entropy.
The ensemble size follows standard deep-ensemble practice ($M{=}5$,
\cite{lakshminarayanan2017simple}); a sensitivity analysis over
$M \in \{1,2,5,10,20\}$ shows the resulting categorization and remedy shares
are stable for $M \ge 5$ (category fractions move $\le 2.6$\,pp between $M{=}5$
and $M{=}20$; \appref{app:msens}), and \appref{app:nontree}
re-derives the full decomposition with bagged non-tree ensembles (MLPs,
logistic regression), with the headline unchanged.

The decomposition separates two distinct sources of error. Aleatoric uncertainty
captures the irreducibility of the data-generating process: it is high in regions
of genuine class overlap. Epistemic uncertainty captures the limits of the current
sample: it is high where the model has seen insufficient training data. This
distinction matters for rebalancing. Boundary-overlap errors (high aleatoric) are
not correctable by adding more data of the existing class; sparse-region errors
(high epistemic) are exactly where additional same-class evidence helps.

By contrast, the k-nearest-neighbour categorizations of Napiera{\l}a \&
Stefanowski \cite{napierala2016types} and related work \cite{smith2014instance}
measure only local same-class geometry, which conflates these two sources: an
instance with $0$ same-class neighbours might be a labelling error, a
Bayes-boundary point in a high-overlap region, or a genuinely under-represented
sparse instance. We use the aleatoric--epistemic decomposition purely as a
\emph{diagnostic} (to localize the boundary/overlap region on real data and
characterize the error population, \S\ref{sec:results}), not as a new
oversampler.

\paragraph{Out-of-bag estimation (no in-sample optimism).} Every
\emph{model-derived} signal that drives the categorization is computed from
\emph{out-of-bag} (OOB) predictions, i.e.\ for each instance, only from the
trees that did not bootstrap it. An instance is an \emph{error} if the OOB
ensemble majority vote disagrees with its label; the true-class probability (TCP,
\S\ref{sec:method:triage}) is the OOB estimate; the noise signals use OOB
confidence and consensus. The remaining signals (the local class ratio) are
properties of the data geometry, not model predictions. The error set and its
categories are therefore \emph{not} in-sample artefacts of forests fit and scored
on the same data. One quantity is scoped differently: the aleatoric/epistemic
values of Eq.~\eqref{eq:uncertainty} are computed over all trees of the fitted
instrument (in-bag included); they do not enter the category rules of
\S\ref{sec:method:triage}, serving as interpretation and validation, and the
strict check below recomputes them fully held-out. As a strict check we additionally re-derive the categories under
inner cross-validation (forests fit on inner-training folds categorize held-out
instances, so even the aleatoric/epistemic decomposition is out-of-fold); the
category distribution and the aleatoric/epistemic separation are essentially
unchanged (\appref{app:oof}). We then categorize each error instance.

\subsection{Three-way error categorization}\label{sec:method:triage}

For each error instance $\mathbf{x}_i$ we compute a \emph{local class ratio}: the
fraction of $k$-nearest neighbours sharing its label, normalized by the global
class frequency. Low values indicate the class is locally under-represented (a
data gap); high values indicate adequate local representation despite the error
(genuine overlap). The categorization rules are:
\begin{description}
\item[\textbf{Cat\,1 (Noise):}] $\mathrm{TCP}(\mathbf{x}_i) < \tau_{\text{noise}}$,
  where TCP is the ensemble true-class probability. Very low TCP indicates likely
  mislabeling; $\tau_{\text{noise}}=0.12$ gates noise from sparse-boundary errors.
\item[\textbf{Cat\,2 (Data-limited):}] not Cat\,1, and
  $\mathrm{class\_ratio}(\mathbf{x}_i) < \tau_{\text{cat2}}$ ($\tau_{\text{cat2}}=0.4$).
  The class is locally under-represented; the error arises from insufficient
  same-class data, not overlap.
\item[\textbf{Cat\,3 (Irreducible):}] neither Cat\,1 nor Cat\,2. The class is
  adequately represented locally yet the model errs, indicating genuine class
  overlap in a model-estimated Bayes-boundary/overlap region. (We reserve the term
  ``Bayes boundary'' for the controlled experiments of \S\ref{sec:results:mechanism}
  where it is known; on real data Cat\,3 is an \emph{estimate} of that region.)
\end{description}
The rules are deliberately operational --- ensemble confidence (TCP) and local
representation --- rather than thresholds on the uncertainty values of
Eq.~\eqref{eq:uncertainty} themselves; the aleatoric--epistemic decomposition
supplies the \emph{interpretation} of the resulting categories, and the held-out
validation confirms it (Cat\,3 aleatoric uncertainty exceeds Cat\,2's on $80\%$
of datasets, \appref{app:oof}).

To remove the imbalance bias of an off-the-shelf ensemble (a majority-biased
ensemble suppresses minority TCP and over-flags minority instances as noise), the
triage ensemble is trained with class-balanced instance weights (inverse class
frequency). Forest probability estimates respond only weakly to such weights
(\S\ref{sec:results:frontier}), and the \emph{unweighted} bagged instruments of
\appref{app:nontree} reproduce the decomposition, so the headline does not hinge
on this choice.

The two thresholds are fixed at $\tau_{\text{noise}}=0.12$, $\tau_{\text{cat2}}=0.4$
across \emph{all} datasets; they are defaults, not tuned hyperparameters. The
categorization is robust to their exact values: sweeping
$\tau_{\text{noise}}\in[0.05,0.30]$ and $\tau_{\text{cat2}}\in[0.2,0.6]$, $74\%$ of
datasets show $<0.5\%$ downstream-accuracy variation across all configurations, and
the category assignments preserve their relative structure throughout
(\appref{app:sensitivity}).

\subsection{Remedy decomposition of the minority deficit}\label{sec:method:remedy}

The triage characterizes \emph{why} an instance is hard; our central instrument
turns this into \emph{which remedy applies} to the instances that actually drive
the deficit --- the misclassified minority. Two instrument choices matter here
and we state them explicitly. First, the deficit being decomposed is that of the
\emph{default, off-the-shelf} model, so the decomposition is measured on the
\emph{unweighted} ensemble; the class-weighted variant of
\S\ref{sec:method:triage} partially enacts the operating-point shift inside the
instrument, leaving a harder residue, and is reported as a sensitivity
(\S\ref{sec:results:remedy}). Second, the error set is the ensemble's
\emph{actual} default prediction: a minority instance is a default error iff the
out-of-bag argmax over its mean probability vector differs from its label (on
binary tasks this coincides with $\Pr(\text{minority})<\tfrac12$; on multiclass
tasks the two differ, and the argmax definition is the valid one). For each such
error we ask, in order:
\begin{description}
\item[\textbf{Threshold-recoverable:}] is it classified \emph{correctly} under
  the decision-rule intervention that predicts the focal minority class whenever
  $\Pr(\text{minority})\ge\tau^\star$ and otherwise the highest-scoring remaining
  class (a one-vs-rest threshold move, valid for binary and multiclass alike),
  where $\tau^\star$ maximizes minority-vs-rest balanced accuracy? The threshold
  is \emph{cross-fitted}: instances are split into five stratified folds, and
  each instance's recovery is judged with a $\tau^\star$ selected on the other
  folds only, so no recovery decision uses a threshold chosen on that instance.
  If recovered, the model's score already places it on the minority side of the
  tuned operating point and only the default cutoff was wrong: the remedy is to
  change the decision rule, no data and no refitting required. This axis, which
  conditions on score-recoverability of an already-misclassified instance, is
  what distinguishes our decomposition from difficulty taxonomies.
\item[\textbf{Data-reducible (epistemic):}] otherwise, if the instance is Cat\,2
  (locally under-represented, high epistemic uncertainty), the remedy is more
  same-class data there (collection or targeted generation).
\item[\textbf{Irreducible (aleatoric):}] otherwise, if it is Cat\,3
  (model-estimated overlap, high aleatoric uncertainty), no remedy in our menu
  recovers it. On real data this is an \emph{estimated}-irreducible, aleatoric
  residue, not a proven Bayes-irreducible one.
\end{description}
(Cat\,1 noise is reported separately.) The $\Pr(\text{minority})$ scores used
here are out-of-bag with respect to model fitting, and $\tau^\star$ is
cross-fitted as above, so the decomposition is held out from both model fitting
and threshold selection (\appref{app:oof}). This three-way, remedy-mapped split
of the \emph{minority error population} is the paper's headline diagnostic; its
empirical shape (\S\ref{sec:results:remedy}) explains why a threshold move recovers
most of the deficit and identifies the minority of cases where data genuinely helps.

\paragraph{Remedy-irreducible is not the Cat\,3 set.} The remedy decomposition is
not a relabelling of the triage, and the \emph{irreducible} remedy class is much
smaller than the Cat\,3 (boundary/overlap) category (Figure~\ref{fig:remedy}). Two
things separate them.
First, the remedy decomposition runs only on the \emph{misclassified minority},
whereas the Cat\,3 fraction is reported over \emph{all} errors. Second, and more
importantly, it gives threshold-recoverability \emph{priority} over the triage
category: an error the triage places in the boundary/overlap region (Cat\,3) is
still counted \emph{threshold-recoverable} whenever the model ranks it above
$\tau^\star$: the default $\tfrac12$ cutoff was wrong, but the \emph{ranking} is
right. The irreducible remedy class therefore contains only the Cat\,3 minority
errors that remain misclassified \emph{after} the threshold move, a strict subset
of Cat\,3. This is why the irreducible remedy fraction ($\sim$17\% of the minority
deficit, \S\ref{sec:results:remedy}) is far below the Cat\,3 share of all errors
($\sim$85\%, \S\ref{sec:results:interventional}): most boundary errors are scored
well enough to be recovered by moving the decision, not by adding or accepting data.

\begin{figure}[t]
\centering
\begin{tikzpicture}[
  font=\small,
  box/.style={draw, rounded corners, align=center, inner sep=3pt, text width=23mm, minimum height=9mm},
  rem/.style={box, fill=blue!8},
  cat/.style={box, fill=black!6},
  q/.style={draw, diamond, aspect=2.4, align=center, inner sep=0pt, text width=17mm, fill=yellow!18},
  lbl/.style={font=\footnotesize\itshape},
  >={Stealth[length=2mm]}
]
\node[box, fill=black!4, text width=46mm] (err)
  {Misclassified minority error\\[1pt]{\footnotesize Cat\,3 $\approx 85\%$ of \emph{all} errors}};
\node[q, below=6mm of err] (q1) {scores above $\tau^\star$?};
\node[rem, left=16mm of q1] (tr) {\textbf{Threshold-}\\\textbf{recoverable}\\$68\%$};
\node[q, below=9mm of q1] (q2) {triage\\category};
\node[cat, below=12mm of q2] (c3) {\textbf{Irreducible}\\Cat\,3, aleatoric\\$17\%$};
\node[cat, left=7mm of c3] (c2) {\textbf{Data-reducible}\\Cat\,2, epistemic\\$6\%$};
\node[cat, right=7mm of c3] (c1) {Noise\\Cat\,1\\$9\%$};
\draw[->] (err) -- (q1);
\draw[->] (q1) -- node[above,lbl]{yes} (tr);
\draw[->] (q1) -- node[right,lbl]{no} (q2);
\draw[->] (q2.south) -- (c3.north);
\draw[->] (q2.south) -- (c2.north);
\draw[->] (q2.south) -- (c1.north);
\end{tikzpicture}
\caption{How the misclassified minority filters into the three remedy classes. The
decomposition asks threshold-recoverability \emph{first}: an error the model already
scores above the cross-fitted threshold $\tau^\star$ is \emph{threshold-recoverable}
whatever its triage category. Only the residue is split by triage into data-reducible
(Cat\,2), irreducible (Cat\,3), or noise (Cat\,1). Because most boundary (Cat\,3)
errors are well-scored, the \emph{irreducible} remedy class ($17\%$) is a strict
subset of Cat\,3 (which is itself $\sim$85\% of \emph{all} errors, \S\ref{sec:results:interventional}), which is why
$17\% \ll 85\%$. Percentages are means over datasets (Table~\ref{tab:reducibility}).}
\label{fig:remedy}
\end{figure}

\begin{table}[t]
\centering
\caption{The triage and remedy-decomposition procedure in pseudocode. All
model-derived quantities are out-of-bag (OOB) or out-of-fold (OOF); nothing is
fit and scored on the same data.}
\label{tab:algorithm}
\footnotesize
\begin{tabular}{@{}rp{0.86\linewidth}@{}}
\toprule
\multicolumn{2}{@{}l}{\textbf{Input:} training split $\mathcal{D}=\{(\mathbf{x}_i,y_i)\}$;
  $\tau_{\text{noise}}{=}0.12$, $\tau_{\text{cat2}}{=}0.4$; $M{=}5$, $T{=}100$, $k$} \\
\multicolumn{2}{@{}l}{\textbf{Output:} triage categories; remedy shares of the minority deficit} \\
\midrule
1 & Fit $M$ random forests ($T$ trees each, class-balanced instance weights); for each
    $\mathbf{x}_i$ collect the OOB tree probabilities $\hat{p}_{m,t}(\mathbf{x}_i)$
    and their mean $\bar{p}(\mathbf{x}_i)$. \\
2 & Error set $E \leftarrow \{i : \text{OOB majority vote} \ne y_i\}$. \\
3 & For each $i \in E$: $\mathrm{TCP}_i \leftarrow \bar{p}(\mathbf{x}_i)[y_i]$;
    aleatoric/epistemic uncertainty via Eq.~(\ref{eq:uncertainty});
    local class ratio $r_i \leftarrow$ fraction of $k$-nearest neighbours sharing
    $y_i$, normalized by the global frequency of $y_i$. \\
4 & \textbf{Triage:} Cat\,1 (noise) if $\mathrm{TCP}_i < \tau_{\text{noise}}$;
    else Cat\,2 (data-limited) if $r_i < \tau_{\text{cat2}}$;
    else Cat\,3 (boundary/overlap). \\
5 & \textbf{Threshold:} $\tau^\star \leftarrow \arg\max_{\tau}$ balanced accuracy
    of the rule $\bar{p}(\text{minority}\mid\mathbf{x}) \ge \tau$ on OOF
    predictions of the training split. \\
6 & \textbf{Remedy decomposition:} for each \emph{minority} $i \in E$:
    \emph{threshold-recoverable} if
    $\bar{p}(\text{minority}\mid\mathbf{x}_i) \ge \tau^\star$;
    else \emph{data-reducible} if Cat\,2; \emph{irreducible} if Cat\,3;
    \emph{noise} if Cat\,1. \\
7 & Return category fractions over $E$ and remedy shares over the minority errors. \\
\bottomrule
\end{tabular}
\end{table}

\subsection{Why boundary interpolation crosses the posterior boundary}\label{sec:method:theory}

The mechanism has a simple analytic basis. Let class~$1$ be the minority, write the
class-posterior $\eta(\mathbf{x})=\Pr(y{=}1\mid\mathbf{x})$, and let the Bayes
boundary be $\mathcal{B}=\{\mathbf{x}:\eta(\mathbf{x})=\tfrac12\}$ with
majority-posterior region $\mathcal{M}=\{\eta<\tfrac12\}$ (where the Bayes-optimal
label is the majority). SMOTE forms a synthetic point
$\mathbf{z}=\mathbf{x}_i+\lambda(\mathbf{x}_j-\mathbf{x}_i)$ with
$\lambda\sim\mathcal{U}(0,1)$, $\mathbf{x}_i$ a minority seed and $\mathbf{x}_j$
one of its minority $k$-nearest neighbours, and labels $\mathbf{z}$ as
\emph{minority}.

\begin{proposition}\label{prop:cross}
Assume $\eta$ is continuous along the segment $[\mathbf{x}_i,\mathbf{x}_j]$.
\begin{enumerate}
\item[(i)] Suppose a minority seed lies in the overlap region,
  $\eta(\mathbf{x}_i)<\tfrac12$, and is interpolated with a neighbour
  $\eta(\mathbf{x}_j)<\tfrac12$. By continuity $\mathcal{M}$ contains an
  $\varepsilon$-neighbourhood of each endpoint, so every \emph{sufficiently short}
  interpolation step from either endpoint yields $\mathbf{z}\in\mathcal{M}$; if in
  addition $\eta$ is quasiconvex along the segment (in particular, approximately
  affine, as in (ii)), then \emph{every} $\mathbf{z}\in[\mathbf{x}_i,\mathbf{x}_j]$
  lies in $\mathcal{M}$. Such a $\mathbf{z}$ carries a minority label in a region
  the Bayes rule assigns to the majority, i.e.\ it acts as label noise.
\item[(ii)] If the segment crosses $\mathcal{B}$, the probability that
  $\mathbf{z}\in\mathcal{M}$ equals
  $\big|\{\lambda\in[0,1]:\eta(\mathbf{x}_i+\lambda(\mathbf{x}_j-\mathbf{x}_i))<\tfrac12\}\big|\in(0,1)$.
  When $\eta$ is approximately affine over the segment with
  $\eta(\mathbf{x}_i)\ge\tfrac12>\eta(\mathbf{x}_j)$, this probability is
  $\dfrac{\tfrac12-\eta(\mathbf{x}_j)}{\eta(\mathbf{x}_i)-\eta(\mathbf{x}_j)}$.
\item[(iii)] Consequently the expected fraction of synthetic minority mass placed
  in $\mathcal{M}$ is non-decreasing in the local overlap (the minority probability
  mass near and beyond $\mathcal{B}$) and in the proportion of seeds drawn from the
  boundary/overlap region.
\end{enumerate}
\end{proposition}

\noindent\emph{Proof sketch.} (i) $\mathcal{M}=\{\eta<\tfrac12\}$ is the open
sublevel set of a continuous function, so it contains an $\varepsilon$-ball around
each interior endpoint and sufficiently short interpolation steps stay in
$\mathcal{M}$. A sublevel set need not be convex, so the segment is guaranteed to
remain in $\mathcal{M}$ \emph{in its entirety} only under the added quasiconvexity
(or approximately-affine) hypothesis; absent it, an interior excursion above
$\tfrac12$ is possible but does not affect the label-noise conclusion for the
points that do fall in $\mathcal{M}$. (ii) follows from the
intermediate value theorem and, in the affine case, from solving
$\eta(\mathbf{x}_i)+\lambda[\eta(\mathbf{x}_j)-\eta(\mathbf{x}_i)]=\tfrac12$ for the
crossing $\lambda^\star$ and integrating the uniform density over
$\{\eta<\tfrac12\}$. (iii) follows by taking expectations over the seed/neighbour
sampling distribution, whose support shifts toward $\mathcal{B}$ as overlap and
boundary-seed proportion grow. \hfill$\square$

This explains three observations below: standard SMOTE places an
\emph{overlap-dependent} fraction of synthetic points across a \emph{known}
boundary (\S\ref{sec:results:mechanism}); augmenting estimated-boundary (Cat\,3)
seeds reproduces the tradeoff while non-boundary seeds do not
(\S\ref{sec:results:interventional}); and boundary-\emph{targeting} variants
(Borderline-SMOTE, ADASYN) raise exactly the boundary-seed proportion of
clause~(iii), so they do not resolve the tradeoff.

\subsection{Strategy menu}\label{sec:method:menu}

We evaluate a menu of 16 rebalancing strategies grouped into five families:
(i)~\emph{do nothing}: train on the raw data;
(ii)~\emph{threshold-moving}: choose the decision threshold on out-of-fold
training predictions to maximize balanced accuracy (binary tasks), leaving the
data untouched;
(iii)~\emph{cost-sensitive}: class-balanced sample weights
\cite{elkan2001foundations}, and triage-informed weighting (upweighting the
correct same-class neighbours of Cat\,2 errors);
(iv)~\emph{generative oversampling}: SMOTE \cite{chawla2002smote} and eight
variants (Borderline-SMOTE, ADASYN, Safe-Level-SMOTE, ProWSyn, MWMOTE,
polynom-fit, Napierala-guided), plus \emph{clean-masked SMOTE} (SMOTE seeded only
from non-boundary, non-noise minority instances, i.e.\ excluding Cat\,1/Cat\,3);
at review we additionally evaluate the modern geometric variant G-SMOTE
\cite{douzas2019geometric} under the same protocol, reported alongside the
pre-registered menu rather than inside it
(\S\ref{sec:results:mechanism}, \S\ref{sec:results:frontier}, \appref{app:oversamplers});
(v)~\emph{balanced ensembles}: balanced random forest \cite{chen2004using} and
EasyEnsemble \cite{liu2009exploratory}.
The two triage-derived members (clean-masked SMOTE, triage weighting)
operationalize the diagnostic; as \S\ref{sec:results} shows they do not beat
the strong non-generative families, and we include them for completeness rather
than as headline methods.

\subsection{Meta-features and prescriptive selection}\label{sec:method:selection}

To ask whether a dataset's structure determines its best strategy, we describe
each dataset by meta-features in three groups: (a)~\emph{error-structure} from the
triage (Cat\,1/2/3 fractions and masses, minority error rate, learnable-minority
fraction); (b)~\emph{complexity / separability}: Fisher's maximum discriminant
ratio (F1) and the boundary/overlap measures N1, N2, N3 \cite{komorniczak2022data};
and (c)~\emph{basic} descriptors (imbalance ratio, size, dimensionality, class
count). A leave-one-dataset-out selector predicts, from these meta-features, which
menu strategy maximizes the target metric on the held-out dataset. We evaluate two
designs: a per-method regressor (predict each strategy's score, take the
argmax) and a direct method-classifier. We compare both against fixed baselines (always-SMOTE,
the best single fixed strategy), an imbalance-ratio-only rule, and a per-dataset
oracle. Significance is assessed by paired Wilcoxon tests across datasets with
bootstrap confidence intervals, and selector stability is checked across random
seeds.
